\documentclass[11pt]{article}
\usepackage[T1]{fontenc}
\usepackage{amsmath,amssymb,amsthm}
\usepackage[margin=1in]{geometry}
\usepackage[hidelinks]{hyperref}

\hypersetup{
  pdftitle={A Uniform Four-Attack Obstruction for the Order Gadget at M = 4 mod 8},
  pdfauthor={Xiaojian Zeng}
}

\newtheorem{theorem}{Theorem}
\newtheorem{lemma}[theorem]{Lemma}

\newcommand{\IM}{I_M}
\newcommand{\OG}{\operatorname{OG}}
\newcommand{\Cfour}{\mathrm{C4}}

\title{A Uniform Four-Attack Obstruction for the Order Gadget\\
at \texorpdfstring{$M\equiv4\pmod 8$}{M = 4 mod 8}}
\author{Xiaojian Zeng\\
PBC School of Finance, Tsinghua University, 100190\\
\texttt{xiaojianzeng@foxmail.com}}
\date{}

\begin{document}
\maketitle

\begin{abstract}
Let $\OG(M)$ be Kasel's finite order gadget on $(M,2M]$: its order must
avoid monotone three-term arithmetic progressions and satisfy fifteen
precedence constraints induced by the attacking values $15$ and $16$.
We prove that $\OG(M)$ is infeasible for every $M\geq20$ with
$M\equiv4\pmod 8$.  In fact, four of its precedence constraints already
suffice.  The proof extends the mirror-flood argument from steps $2$ and $4$
to every positive even step and uses step $8$ after splitting the target
class modulo $16$.  Exact finite checks audit the displayed proof schemas but
are not used in the universal argument.  The result concerns this finite
gadget and residue class; it does not decide the Erd\H{o}s--Graham two-set
permutation problem.
\end{abstract}

\section{Statement and context}

For an integer $M\geq20$, put
\[
  \IM=(M,2M]\cap\mathbb Z,
  \qquad b_j=M+j,
  \qquad t_i=2M-i.
\]
A strict total order $\prec$ on $\IM$ is \emph{AP-free} if no numerical
three-term arithmetic progression $a<b<c$ in $\IM$ is monotone in the order:
neither $a\prec b\prec c$ nor $c\prec b\prec a$ may hold.  Equivalently, the
numerical midpoint is before both endpoints or after both endpoints.

The order gadget $\OG(M)$ was introduced in Definition~15 of
Kasel~\cite{Kasel2026}.  Besides AP-freeness, it imposes the fifteen
precedence relations obtained from the attacking values $15$ and $16$.
The source proves a three-relation obstruction when $M\equiv0\pmod8$, and it
reports finite solver observations for other attack cores, but it does not
give a uniform theorem for the present residue class
\cite[Proposition~19, Theorem~22, and Remark~30]{Kasel2026}.  The underlying
two-set question arose in work of Davis, Entringer, Graham, and
Simmons~\cite{DavisEtAl1977} and was subsequently emphasized by Erd\H{o}s and
Graham~\cite{ErdosGraham1979}.

Consider the four constraints
\begin{equation}\label{eq:C4}
\begin{aligned}
 P&:\ t_{13}\prec b_1, &
 Q&:\ t_{11}\prec b_2,\\
 A_1&:\ t_5\prec b_5, &
 A_3&:\ t_{10}\prec b_3.
\end{aligned}
\end{equation}
They are all among the fifteen constraints of $\OG(M)$.

\begin{theorem}\label{thm:main}
For every integer $M\geq20$ with $M\equiv4\pmod8$, no AP-free strict total
order on $\IM$ satisfies the four constraints in \eqref{eq:C4}.  Consequently
$\OG(M)$ is infeasible for every such $M$.
\end{theorem}

The theorem is a uniform statement, not an extrapolation from finite solver
outcomes.  Its proof occupies Sections~\ref{sec:tools}--\ref{sec:proof}.

\section{Ladders and a phase-blind flood lemma}\label{sec:tools}

For a numerical progression $a<b<c$, midpoint extremality gives the four
implications
\begin{equation}\label{eq:R}
\begin{aligned}
 a\prec b&\Longrightarrow c\prec b, &
 c\prec b&\Longrightarrow a\prec b,\\
 b\prec c&\Longrightarrow b\prec a, &
 b\prec a&\Longrightarrow b\prec c.
\end{aligned}
\end{equation}
Fix a positive step $g$.  On the elements of $\IM$ in one residue class
modulo $g$, consecutive numerical triples and \eqref{eq:R} force a zigzag:
one of the two classes modulo $2g$ precedes both neighboring rungs, and the
other class follows both neighboring rungs.  We call the former rungs
\emph{leaders} and the latter \emph{trailers}.  Either zigzag phase may occur.

We use the following phase-blind form of the flood lemma.  For $g=2,4$ this
is the mechanism of Kasel's Lemma~26~\cite{Kasel2026}; the same elementary
induction applies to every positive even $g$.

\begin{lemma}[general flood lemma]\label{lem:flood}
Let $g$ be a positive even integer, and let $C$ be the intersection of
$\IM$ with one residue class modulo $g$.  Suppose that $c\in\IM$ is displaced
by $g/2$ from that residue class.  Define
\[
 W_c=\{w\in C:2c-w\in\IM\}.
\]
If $c\prec w_0$ for one $w_0\in W_c$, then $c\prec w$ for every $w\in W_c$.
If $w_0\prec c$ for one $w_0\in W_c$, then $w\prec c$ for every $w\in W_c$.
The conclusion is independent of the zigzag phase of the $g$-ladder.
\end{lemma}

\begin{proof}
The two elements $c-e$ and $c+e$ of any admissible mirror pair belong to
opposite classes modulo $2g$, because $e\equiv g/2\pmod g$.  Thus exactly one
is a leader.  The progression $(c-e,c,c+e)$ and \eqref{eq:R} first copy any
seed comparison with $c$ to the other member of its mirror pair.

Suppose first that the pair at distance $e$ lies above $c$ in the order.  To
reach an admissible pair at distance $e+g$, choose the leader in the first
pair.  It precedes its outward ladder neighbor, so transitivity places that
neighbor above $c$; the mirror rule then reaches the other member of the new
pair.  To reach an admissible pair at distance $e-g$, choose the trailer in
the nearer pair.  Its outward neighbor at distance $e$ is a leader, so the
known comparison, the leader--trailer edge, and transitivity reach the nearer
trailer; reflection reaches its partner.

For comparisons directed toward $c$, use the same two moves with inequalities
reversed: at distance $e+g$ its leader precedes the already-reached inward
neighbor, while at distance $e-g$ the nearer leader precedes its
already-reached outward neighbor.  Reflection completes each pair.  The
admissible positive distances form a finite consecutive interval in the
class $g/2\pmod g$.  Induction in both directions therefore reaches every
member of $W_c$, in either zigzag phase.
\end{proof}

We shall also repeatedly use a direct corollary.  On the odd step-$2$ ladder,
any odd value $c$ is displaced by $2$ from the opposite class modulo $4$.
An adjacent leader--trailer comparison seeds Lemma~\ref{lem:flood} with
$g=4$.  Thus, if $c$ is a leader, it precedes every mirror-admissible trailer
in the opposite class; if $c$ is a trailer, every mirror-admissible leader in
the opposite class precedes $c$.

\section{The four-attack obstruction}\label{sec:proof}

Throughout this section assume $M\geq20$ and $M\equiv4\pmod8$, and put
\[
  m=\frac{3M}{2}.
\]
Then $m\equiv2\pmod4$ and $m\in\IM$.

\begin{lemma}\label{lem:phase}
In an AP-free order, $Q$ and $A_3$ imply $b_5\prec b_3$.
\end{lemma}

\begin{proof}
Suppose instead that $b_3\prec b_5$.  The odd step-$2$ ladder then has leader
class $3\pmod4$.  Set
\[
  c^*=m-1\equiv1\pmod4,
  \qquad
  c^{**}=m+1\equiv3\pmod4.
\]
Compare $t_5$ and $m$ in the total order.

If $t_5\prec m$, the step-$2$ flood centered at $m$, seeded by the mirror
pair $\{b_5,t_5\}$, puts every odd value below $m$.  In particular
$c^*\prec m$.  The odd-ladder phase and the step-$4$ flood centered at $c^*$
give $b_3\prec c^*$; the reflection of $b_3$ about $c^*$ is $t_5$.
The step-$2$ flood centered at $c^*$, seeded by $c^*\prec m$, then gives
$c^*\prec t_{10}$; the reflection of $t_{10}$ is $b_8$.  Constraint $A_3$
produces the strict cycle
\[
  b_3\prec c^*\prec t_{10}\prec b_3.
\]

If $m\prec t_5$, the same polar step-$2$ flood puts $m$ below every odd
value, so $m\prec c^{**}$.  Since $c^{**}$ is an odd leader, the step-$4$
flood at $c^{**}$ gives $c^{**}\prec t_{11}$, whose reflection is $b_{13}$.
The step-$2$ flood at $c^{**}$, seeded by $m\prec c^{**}$, gives
$b_2\prec c^{**}$, whose reflection is $t_0$.  Constraint $Q$ now gives
\[
  c^{**}\prec t_{11}\prec b_2\prec c^{**}.
\]
Both cases are impossible.  All displayed points and reflections lie in
$\IM$ for $M\geq20$.
\end{proof}

\begin{lemma}\label{lem:flip}
In an AP-free order satisfying $b_5\prec b_3$, the constraints $P$ and $A_1$
cannot both hold.
\end{lemma}

\begin{proof}
The assumed orientation makes the odd leader class $1\pmod4$.

First suppose $M\equiv4\pmod{16}$.  Then $m\equiv6\pmod8$.  Put
\[
  c=m-5\equiv1\pmod8,
  \qquad
  d=m-1\equiv5\pmod8.
\]
Both $c$ and $d$ are odd leaders.  A step-$4$ flood at $c$ gives
$c\prec t_{13}$, with reflection $b_3$; then $P$ gives $c\prec b_1$.
The step-$8$ flood centered at $c$, over the class of $b_1$, gives
$c\prec d$.  Its seed reflection is $t_{11}$ and its target reflection is
$m-9$.

Similarly, a step-$4$ flood at $d$ gives $d\prec t_5$, with reflection
$b_3$, and $A_1$ gives $d\prec b_5$.  At $M=20$ one has $b_5=c$, already
contradicting $c\prec d$.  For every remaining value in this subclass,
$M\geq36$.  The step-$8$ flood centered at $d$, over the class of $b_5$,
then gives $d\prec c$; its seed reflection is $t_7$ and its target reflection
is $m+3$.  This contradicts $c\prec d$.  The tight boundary for the first
step-$8$ flood is $m-9\geq M+1$, which is equivalent to $M\geq20$.

Now suppose $M\equiv12\pmod{16}$.  Then $m\equiv2\pmod8$.  Put
\[
  c=m-1\equiv1\pmod8,
  \qquad
  d=m-5\equiv5\pmod8.
\]
A step-$4$ flood at $c$ gives $c\prec t_{13}$, with reflection $b_{11}$,
and $P$ gives $c\prec b_1$.  The step-$8$ flood at $c$ over the class of
$b_1$ gives $c\prec d$; its seed reflection is $t_3$ and its target
reflection is $m+3$.

For the reverse comparison, a step-$4$ flood at $d$ gives
$d\prec t_{13}$, with reflection $b_3$, and hence $d\prec b_1$ by $P$.
Numerically $(b_1,d,t_{11})$ is an arithmetic progression with midpoint $d$,
so \eqref{eq:R} yields $d\prec t_{11}$.  The value $t_{11}$ is an odd leader;
a step-$4$ flood centered there gives $t_{11}\prec t_5$, with reflection
$t_{17}$.  By $A_1$, we obtain $d\prec b_5$.  Finally the step-$8$ flood at
$d$, over the class of $b_5$, gives $d\prec c$; its seed reflection is
$t_{15}$ and its target reflection is $m-9$.  This contradicts $c\prec d$.
The least value in this subclass is $M=28$, where all named values and
reflections already lie in $\IM$.

The two cases exhaust $M\equiv4\pmod8$.
\end{proof}

\begin{proof}[Proof of Theorem~\ref{thm:main}]
Constraints $Q$ and $A_3$ force $b_5\prec b_3$ by Lemma~\ref{lem:phase}.
Lemma~\ref{lem:flip} says that this orientation is incompatible with $P$ and
$A_1$.  Thus the subsystem \eqref{eq:C4} is infeasible.  Since the full
gadget contains these four constraints, the same is true of $\OG(M)$.
\end{proof}

\section{Computational checks, scope, and assistance}

The written schemas were audited by exact computation.  One checker verifies
every stated congruence, arithmetic progression, reflection, block boundary,
and transitivity step for all $2{,}046$ target values
$20\leq M\leq16{,}380$.  Separately, a tournament-with-transitivity encoding
and an all-different integer-rank encoding agree on infeasibility at
$M=20,28,36,116,204$.  These finite calculations are corroboration only:
the quantified conclusion follows from the proof above.  This separation
between a theorem's stated scope and exact finite certification is also
emphasized in related finite-certificate work~\cite{ZengEtAl2026}.

The closest inspected source is Kasel's preprint~\cite{Kasel2026}; exact-title
and equivalent-description searches conducted on 7 September 2026 located no
other primary paper addressing this order gadget.  This is a bounded search
assessment, not a claim of priority or of the absence of unindexed work.
The theorem resolves the specified full-gadget assertion on one residue class,
but it does not prove infeasibility in the other residue classes and does not
settle whether the positive integers can be partitioned into two admissibly
permutable sets.

OpenAI Codex assisted with literature search, proof-schema checking, exact
finite-model code, adversarial audit organization, and initial manuscript
preparation.  Language-model output is not used as mathematical evidence.
The universal conclusion rests on the human-readable argument above; the
reported machine checks use exact discrete constraints and no floating-point
tolerances.  The author takes responsibility for the statement and proof.

\bibliographystyle{plain}
\bibliography{references}

\end{document}
